% Generated draft: requires mathematical and source review.
\documentclass[11pt]{article}
\usepackage[margin=25mm]{geometry}
\usepackage{fontspec}
\setmainfont{Latin Modern Roman}
\usepackage{amsmath,amssymb,mathtools,unicode-math}
\setmathfont{Latin Modern Math}
\usepackage{enumitem,needspace,xurl,hyperref}
\setlength{\parindent}{0pt}
\setlength{\parskip}{5pt}
\setlength{\emergencystretch}{3em}
\newcommand{\N}{\mathbb{N}}
\newcommand{\Z}{\mathbb{Z}}
\newcommand{\Q}{\mathbb{Q}}
\newcommand{\R}{\mathbb{R}}
\newcommand{\C}{\mathbb{C}}
\newcommand{\F}{\mathbb{F}}
\newcommand{\A}{\mathbb{A}}
\newcommand{\PP}{\mathbb{P}}
\newcommand{\E}{\mathbb{E}}
\newcommand{\eps}{\varepsilon}
\newcommand{\dd}{\,\mathrm{d}}
\begin{document}
\section*{UnsolvedMath ID: 10400177}
\textbf{Source problem code:} catalogue entry
\subsection{Definitions and mathematical statement}
Turaev asks to relate the surgery formula for the Casson--Walker--Lescop invariant described immediately before this problem to Lescop's formula [251]. If \(D\) is surgery data in a common admissible class presenting a three-manifold \(M_D\), and \(F(D)\) and \(F_{\mathrm{Lescop}}(D)\) denote the two expressions after translating conventions, the requested comparison should explain their relationship and recover the same normalized invariant. The preceding formula, Lescop's formula and their admissible surgery-data classes are absent; \(F\) and \(F_{\mathrm{Lescop}}\) are placeholders and no particular link formula is fabricated.

The supplied normalization footnotes are essential:
\[\lambda_{\mathrm{CW}}(M)=2\lambda_{\mathrm C}(M)
\quad\text{for an integral homology three-sphere }M,\]
\[\lambda_{\mathrm{CWL}}(M)=\frac{|H_1(M;\Z)|}{2}\lambda_{\mathrm{CW}}(M)
\quad\text{for a rational homology three-sphere }M.\]
An integral homology three-sphere has integral homology that of \(S^3\); a rational homology three-sphere has rational homology that of \(S^3\), and then \(|H_1(M;\Z)|\) is the order of its finite first homology group. These formulas imply that the stated conventions give \(\lambda_{\mathrm{CWL}}=\lambda_{\mathrm C}\) on integral homology spheres. Orientation conventions and any wider class of three-manifolds covered by the omitted surgery formulas must still be specified.

The following paragraph, attributed to C. Lescop, supplies context rather than another explicit question: it describes Casson's 1984 integer-valued invariant as an appropriate count of \(SU(2)\)-representations of the fundamental group, citing [3,159], and reports that Cappell, Lee and Miller [76] extended the counting approach to compact Lie groups to obtain further invariants of integral homology spheres. The precise counting and Lie-group hypotheses are not supplied, and this report is not used to invent a second target.
\subsection{Short English statement}
Explain the connection between the preceding surgery formula and Lescop's formula, preserving the stated factors of two and homology-group order. Both formulas themselves are missing from the extracted record.
\subsection{Sources}
\begin{itemize}
\item Local UnsolvedMath registry, permanent ID \(10400177\), exact source code \texttt{catalogue entry}.
\item V. Turaev, attribution in Ohtsuki, \emph{Problems on invariants of knots and 3-manifolds} (2002), Problem 10.2, PDF page 142, normalization footnotes 32--33: \url{https://msp.org/gtm/2002/04/gtm-2002-04-024s.pdf}. Registry access 2026-07-29; pdf-font-structure extraction.
\item Lescop [251] is the comparison reference. The following C. Lescop paragraph cites [3,159] for Casson's construction and Cappell--Lee--Miller [76]. Full entries are not supplied.
\end{itemize}
\end{document}
